\section{Research Experience}

\begin{research}
{Starspot Inference with Light Curve Inversion Techniques}
{American Museum of Natural History}
{Dr. Lucy Lu}
{Aug 2023 -- Sep 2025}
    \item Developed computational approaches for starspot inference using light curve inversion techniques.
    \item Applied time-series analysis and probabilistic modeling to characterize stellar surface features.
    \item Implemented Python-based scientific computing workflows for astrophysical inference and model evaluation.
\end{research}

\begin{research}
{Towards Mapping Brown Dwarf and Giant Exoplanet Atmospheres}
{American Museum of Natural History}
{Dr. Johanna Vos}
{May 2021 -- Aug 2023}
    \item Investigated atmospheric variability in brown dwarfs and giant exoplanets through analysis of photometric light curves.
    \item Developed computational workflows for mapping cloud structures from observational time-series data.
    \item Applied quantitative modeling techniques to infer atmospheric properties from astronomical observations.
\end{research}

\begin{research}
{Chemodynamically Characterizing the Jhelum Stellar Stream with APOGEE-2}
{Sloan Digital Sky Survey Faculty and Student Team}
{Dr. Allyson Sheffield}
{May 2020 -- May 2021}
    \item Analyzed APOGEE-2 spectroscopic data to identify candidate stellar members of the Jhelum stellar stream.
    \item Applied statistical and computational analysis methods to characterize stellar populations.
    \item Contributed to a refereed publication in \textit{The Astrophysical Journal}.
\end{research}

\begin{research}
{Looking at Star Formation Through Chemistry}
{American Museum of Natural History}
{Dr. Allyson Sheffield}
{Aug 2018 -- May 2020}
    \item Analyzed stellar spectra to derive chemical abundances and characterize stellar populations.
    \item Processed and normalized echelle spectra and measured equivalent widths of spectral lines.
    \item Applied quantitative techniques to spectroscopic datasets.
\end{research}

\begin{research}
{Two-Point Statistics in the Star Forming ISM}
{Center for Computational Astrophysics, Flatiron Institute}
{Dr. Chang-Goo Kim}
{Jun 2018 -- Aug 2018}
    \item Performed statistical analysis of metallicity correlations in star-forming regions.
    \item Investigated relationships among stellar metallicity, velocity, and star formation rate.
\end{research}